\section{Обычный обратный поиск и вычислимый сертификат}
\label{sec:model}
Пусть $Q\subset\R^d$ — непустое компактное выпуклое множество диаметра
$D>0$. Пусть $f$ выпукла и принадлежит $C^2$ в окрестности $Q$,
$f_Q^*=\min_Q f$ и $\delta(x)=f(x)-f_Q^*$. Предполагается, что
\begin{equation}\label{eq:curvature}
 \norm{\nabla^2f(z)}_2\le H_0+H_1\delta(z),\qquad z\in Q,
 \quad H_0,H_1\ge0.
\end{equation}
Опорное значение в этом условии — минимум на допустимом множестве.
Ограничение относительно меньшего безусловного инфимума не влечёт
того же ограничения относительно условного минимума без изменения
$H_0$. При $D=0$ единственная допустимая точка решает задачу.

В исходной точке $x_k$ вызывается точный оракул линейной минимизации
(LMO), после чего вычисляются FW-зазор и накопленный сертификат:
\begin{align}
 s_k&\in\arg\min_{v\in Q}\inner{\nabla f(x_k)}{v},\quad
 G_k=\inner{\nabla f(x_k)}{x_k-s_k},\label{eq:gap}\\
 \ell_k&=\max_{0\le i\le k}\{f(x_i)-G_i\},\quad
 U_k=f(x_k)-\ell_k.\label{eq:certificate}
\end{align}
Выпуклость даёт $0\le\delta_k\le U_k\le G_k$. Известную корректную
нижнюю границу можно включить в $\ell_k$, но она не обязательна.
Нулевой зазор удостоверяет оптимальность. Остановка происходит при
$U_k\le\varepsilon$.
Конструкция лучшей накопленной касательной границы известна ранее
\cite[метод 1, соглашение о максимизации]{freundgrigas}; новый анализ
касается её стоимости при зависящей от функции кривизне, а не самого
сертификата.

\paragraph{Алгоритм.}
Выберем $M>0$ и положим $L_{-1}=M$. В каждой нетерминальной исходной
точке зададим $d_k=s_k-x_k$, $q_k=\norm{d_k}$ и начнём с
$L=L_{k-1}/2$. Повторяем
\begin{equation}\label{eq:btstep}
 \alpha=\min\{1,G_k/(Lq_k^2)\},\qquad
 f(x_k+\alpha d_k)\le f(x_k)-\alpha G_k
                         +\tfrac12L\alpha^2q_k^2
\end{equation}
до выполнения неравенства, удваивая $L$ после каждого отказа.
Сохраняем $L_k=L$, обновляем $x_{k+1}=x_k+\alpha d_k$ и повторно
используем принятое значение функции. В ходе поиска заново вычисляются
лишь пробные значения функции. Это стандартная схема квадратичного
обратного поиска~\cite{pedregosa}. При фиксированном $M$ алгоритму не
нужны $D,H_0,H_1$ и $f_Q^*$. Необязательный выбор $M=U_0/D^2$
использует границу диаметра и первые оплаченные вызовы функции,
градиента и LMO.

\begin{lemma}[Локальная приёмка]\label{lem:local}
Для допустимого направления спуска $d$ с
$g=-\inner{\nabla f(x)}d>0$ и верхней границей шага $a_{\max}>0$
проверка~\eqref{eq:btstep} проходит для шага
\[
\alpha=\min\{a_{\max},g/(L\norm d^2)\},
\]
если
$L\ge H_0+H_1\delta(x)$. Любая принятая проба, независимо от этого
достаточного порога, удовлетворяет
\begin{equation}\label{eq:acceptedprogress}
 f(x)-f(x+\alpha d)\ge\tfrac12\alpha g.
\end{equation}
Следовательно, поиск завершается и
\begin{equation}\label{eq:localenvelope}
 L_k\le\max\{L_{k-1}/2,\ 2(H_0+H_1\delta_k)\}.
\end{equation}
Допускается также $a_{\max}>1$.
\end{lemma}
\begin{proof}
Запишем $h(t)=f(x+td)-f(x)$ и $K=H_0+H_1\delta(x)$.
Пока $h\le0$, ограничение гессиана даёт $h''(t)\le K\norm d^2$.
В начальный момент $h'(0)=-g<0$. При $K>0$ первый возврат к нулю
до $\min\{a_{\max},2g/(K\norm d^2)\}$ противоречит
$h(t)\le-gt+K\norm d^2t^2/2<0$. То же интегрирование доказывает
квадратичную верхнюю оценку вплоть до заданной пробы; при $K=0$
получается линейная верхняя модель. Её кривизна не превосходит $L$,
поэтому проверка проходит. Для любой принятой пробы
$L\alpha\norm d^2\le g$, откуда следует~\eqref{eq:acceptedprogress}.
Если был отказ, последняя отклонённая оценка меньше $K$, поэтому
принятая меньше $2K$; иначе она равна $L_{k-1}/2$.
\end{proof}

Таким образом, $f_k$ и $U_k$ не возрастают. В анализе используется
кривизна лишь на отрезке спуска, а не единая верхняя модель во всех
точках произвольного отрезка.

\begin{lemma}[Потенциал запаздывания]\label{lem:lag}
Пусть индексы фаз исходных точек не убывают и принадлежат
$\{0,\ldots,J-1\}$. Пусть положительные верхние границы удовлетворяют
$C_j/2\le C_{j+1}\le C_j$. Предположим, что обновление с $L_k>C_j$
в своей исходной фазе обязательно имеет $L_k=L_{k-1}/2$.
Тогда число таких запаздывающих обновлений не превышает
\[
 \ceil{\log_2\max\{1,M/C_0\}}+J-1.
\]
Утверждение применимо к любому префиксу и допускает пропуск начальных фаз.
\end{lemma}
\begin{proof}
Перед обновлением в фазе $j$ используем неотрицательный целочисленный
потенциал
\[
 p=\max\{0,\ceil{\log_2(L_{k-1}/C_j)}\}.
\]
Запаздывающее обновление уменьшает его на единицу; после незапаздывающего
обновления потенциал перед сменой фазы равен нулю. Переход вперёд
на $m$ фаз увеличивает его не более чем на $m$. От виртуальной
нулевой фазы сумма всех переходов, включая начальный, не превосходит
$J-1$. Неотрицательность даёт требуемый подсчёт.
\end{proof}

\begin{theorem}[Наблюдаемая остановка и полная стоимость функции]\label{thm:fw}
Для $0<\varepsilon<U_0$ положим
\[
 J=\ceil{\log_2(U_0/\varepsilon)},\qquad
 P=\ceil{\log_2\max\{1,MD^2/U_0\}}.
\]
Обычный обратный поиск достигает $U_N\le\varepsilon$ за
\begin{equation}\label{eq:fwcount}
 N\le P+\ceil{(6+16H_1D^2)J+32H_0D^2/\varepsilon}
\end{equation}
обновлений. Он использует $N+1$ градиентов и LMO, включая финальную
проверку. При фактическом $N\ge1$ все вычисления функции, включая
отклонённые пробы, учтены в оценке
\begin{equation}\label{eq:functioncount}
 N_f\le1+2N+\ceil{\log_2\max\{1,2(H_0+H_1U_0)/M\}}.
\end{equation}
При $N=0$ достаточно одного вычисления функции, градиента и LMO.
Оценки предполагают точные оракулы и точную арифметику.
\end{theorem}
\begin{proof}
Положим $b_j=U_0/2^j$ и отнесём нетерминальную исходную точку
к фазе $j$, если $b_j/2<U_k\le b_j$, где $j=0,\ldots,J-1$.
Используем
\[
 C_j=\max\{2(H_0+H_1b_j),\ b_j/(2D^2)\}.
\]
Эти границы убывают не более чем вдвое. По лемме~\ref{lem:local}
принятая оценка выше $C_j$ обязательно получена делением пополам
на первой пробе. Поскольку $C_0\ge U_0/(2D^2)$, лемма~\ref{lem:lag}
ограничивает число запаздывающих обновлений величиной $P+J$.

Для незапаздывающего обновления фазы $j$ неравенства
$G_k\ge U_k>b_j/2$ и $q_k\le D$ дают
\[
 f_k-f_{k+1}\ge p_j:=\min\{b_j/4,\ b_j^2/(8C_jD^2)\}.
\]
Каждое обновление уменьшает $f$, а суммарное уменьшение внутри фазы
не превосходит ошибки по функции в её начале, то есть не более $b_j$.
Поэтому число незапаздывающих обновлений фазы не превосходит
\[
 \lfloor b_j/p_j\rfloor+1
 \le5+16H_1D^2+16H_0D^2/b_j.
\]
Суммирование с использованием $\sum_{j<J}1/b_j<2/\varepsilon$
доказывает~\eqref{eq:fwcount}. Оценка ограничивает любой
незавершённый префикс, что также устанавливает остановку.

Локальная огибающая и монотонность $U$ дают
$L_k\le\max\{M/2^{k+1},2(H_0+H_1U_0)\}$.
Если $r_k$ — число отказов, то $L_k=2^{r_k-1}L_{k-1}$, поэтому
число вычисленных пробных значений в точности равно
\[
 \sum_{k<N}(r_k+1)=2N+\log_2(L_{N-1}/M).
\]
Добавляя начальное значение и используя огибающую, получаем
\eqref{eq:functioncount}. Принятые значения сохраняются.
\end{proof}

Последнее тождество стандартно~\cite{pedregosa}. При $M=U_0/D^2$
имеем $P=0$, а добавка к числу вычислений функции не зависит от
$\varepsilon$. Для порядка оценки эта нормировка необязательна.

\begin{proposition}[Инициализация заметна даже на аффинных функциях]
\label{prop:init}
Пусть $Q=[-1,1]$, $f(x)=-cx$, $c>0$, $x_0=-1$ и $L_{-1}=M>0$.
Для $0<\varepsilon<2c$ обратный поиск с делением пополам использует
ровно
\[
 N=\ceil{\log_2\left(1+\frac Mc(1-\varepsilon/(2c))\right)}
\]
обновлений. Следовательно, уменьшение $U_0$ вдвое может требовать
$\Theta(\log(M/c))$ обновлений.
\end{proposition}
\begin{proof}
Каждая проверка принимает первую пробу: $L_k=M/2^{k+1}$,
$x_{k+1}=\min\{1,x_k+c/L_k\}$. Касательная нижняя оценка всегда
равна $-c$, поэтому $x_N=\min\{1,-1+(2c/M)(2^N-1)\}$ и
$U_N=c(1-x_N)$. Решение неравенства остановки даёт формулу.
\end{proof}
Эта нижняя оценка относится к заданному правилу деления пополам,
а не к произвольным алгоритмам первого порядка.
При $M=U_0/D^2=c/2$ первый шаг решает задачу.

\paragraph{Сравнение с глобальной границей.}
Для $f(x)=\cosh(x)-1$ на $[-R,R]$ имеем $H_0=H_1=1$, тогда как
глобальная константа Липшица градиента равна $L_Q=\cosh R$.
При старте из $R$ получаем $U_0=2R\sinh R$. Для фиксированных
точности и $M$ наша оценка числа обновлений растёт как
$O(R^3+R^2/\varepsilon)$ при $R\to\infty$, в отличие от
экспоненциального слагаемого $L_QD^2/\varepsilon$, получаемого
подстановкой $H_0=L_Q,H_1=0$ в ту же сертификатную теорему.
Здесь сравниваются верхние оценки, а не время работы или оракульная
нижняя граница; точная минимизация на отрезке решает этот конкретный
одномерный пример за один шаг. Другие сравнения шагов приведены в
приложении~\ref{app:explicit}.
